\section{Method}
\subsection{Problem Setup}
For each category $c$, we observe a support set $\mathcal{S}_c=\{x_i\}_{i=1}^{k}$ containing only normal images. The test set contains normal and anomalous images, but anomaly labels are used only for evaluation, not for fitting the main detector or calibrator. The goal is to produce: (i) a raw anomaly ranking score $s(x)$, (ii) a probability-like reliability score $\hat{p}(y=1\mid x)$ or conformal anomaly confidence, and (iii) optional pixel-level maps.

\subsection{Frozen Patch Features and Subspace Ranking}
We extract patch features with a frozen DINOv2 ViT-S/14 backbone. Let $z_{ij}\in\mathbb{R}^{d}$ denote patch $j$ from image $i$. A PCA/subspace model is fit on support patches. With support mean $\mu$ and retained principal directions $U\in\mathbb{R}^{d\times m}$, the patch residual is
\begin{equation}
    r(z)=\left\|(z-\mu)-UU^{\top}(z-\mu)\right\|_2 .
\end{equation}
Patch residuals form an anomaly map. The image-level ranking score is a high-residual aggregation,
\begin{equation}
    s(x)=\operatorname{Agg}_{j}\, r(z_j),
\end{equation}
where $\operatorname{Agg}$ is implemented as a maximum or top-tail statistic. Importantly, this raw score is the quantity used for AUROC and AP. Calibration routes do not change the ranking score.

\subsection{Vector and Shift-Aware Calibration}
The first reliability family is post-hoc logistic calibration. A scalar Platt scaler maps $s(x)$ to a probability, while Vector Platt uses a richer descriptor
\begin{equation}
    \phi(x)=\left[s_{\mathrm{pca}}(x),\; s_{\mathrm{head}}(x),\; |\tilde{s}_{\mathrm{pca}}(x)-\tilde{s}_{\mathrm{head}}(x)|\right],
\end{equation}
where $s_{\mathrm{head}}$ is a lightweight synthetic-anomaly head score when available and tildes denote normalized scores. The calibrated probability is
\begin{equation}
    \hat{p}_{\theta}(x)=\sigma(w^{\top}\phi(x)+b).
\end{equation}
Shift-Aware Platt augments $\phi(x)$ with shift descriptors such as residual mean, residual concentration, and support/test domain statistics. Weighted Platt further reweights samples by density-ratio-style estimates. These variants are useful for structured shifts but can over-adapt, motivating a conservative conformal route.

\subsection{LOIO Conformal Reliability}
For $k\ge2$, CRR constructs normal calibration residuals by leave-one-image-out support splitting. For each support image $x_i$, a subspace is fit on $\mathcal{S}_c\setminus\{x_i\}$ and the held-out image receives a residual score $s_{-i}(x_i)$. The resulting set
\begin{equation}
    \mathcal{R}_{\mathrm{LOIO}}=\{s_{-i}(x_i):x_i\in\mathcal{S}_c\}
\end{equation}
serves as normal nonconformity evidence. For a test image $x$, the conformal p-value is
\begin{equation}
    p_{\mathrm{LOIO}}(x)=\frac{1+\sum_{r\in\mathcal{R}_{\mathrm{LOIO}}}\mathbf{1}\{r\ge s(x)\}}{1+|\mathcal{R}_{\mathrm{LOIO}}|}.
\end{equation}
We report $p_{\mathrm{conf}}(x)=1-p_{\mathrm{LOIO}}(x)$ as a probability-like anomaly reliability view. It should be interpreted as relative extremeness against held-out support normals, not as a fully supervised posterior probability. Weighted conformal replaces the uniform count by density-ratio weights and reports effective sample size; empirically, it is less stable than LOIO on the main benchmark.

\subsection{Reliability Routing and Selective Diagnostics}
CRR evaluates multiple reliability routes: Vector Platt, Shift-Aware Platt, weighted Platt, LOIO conformal, weighted conformal, and simple rule-based gates. The main route is fixed LOIO conformal because it does not require validation-label expert selection and is stable in the full VisA corruption benchmark. Gated variants are used diagnostically to test whether shift-aware experts can improve reliability without harming safe baselines. We additionally compute entropy, expert disagreement, conformal confidence, and effective sample size as selective-risk signals for flagging uncertain decisions.
